\chapter{Conclusions and Outlook}
\label{ch:conclusion}

% --------------------------------------------------------------
\section{Conclusions}
\label{sec:conclusion-achievements}
% --------------------------------------------------------------

This thesis set out to build the layers needed to run a \textsc{FeynRules}-class
derivation of Feynman rules inside the Python ecosystem, and to validate the
tensor expressions it produces reliably. Both aims have been achieved.

\textsc{FeynPy} allows a model to be defined from Python objects describing
indices, fields, gauge groups and parameters, together with a Lagrangian
written in a notation close to the one used on paper. From this information it
derives tree-level Feynman rules as \textsc{Symbolica} tensor expressions. The
implementation is generic rather than tied to a particular model: covariant
derivatives, field strengths, gauge fixing and ghost sectors are constructed
from the model declarations themselves. In the QCD example, four written
Lagrangian terms generate the expected two-point terms together with the
quark--gluon, ghost--gluon, three-gluon and four-gluon interactions.

A compiled Lagrangian can also be manipulated after construction. Field
transformations are used, for example, to move from gauge to physical bases,
while gauge and BRST variations can be evaluated directly. These operations
use the same underlying representation rather than separate model-specific
implementations.

The main result of the thesis is the validation of the generated vertices.
Comparing two independently derived tensor expressions requires more than
comparing their literal form: dummy indices may be renamed, factors reordered,
and equivalent tensor structures written differently. The canonicalisation
procedure of \Cref{sec:validation-canonical} provides a common representation
in which these expressions can be compared, so equality is tested as a tensor
identity rather than as a string match.

With this procedure, the packaged Standard Model agrees with its
\textsc{FeynRules} reference at $163/163$ vertices, the unbroken
background-field Standard Model at $67/67$, and the Green-basis SMEFT at
$184/184$ reference lines. Of the latter, $161$ agree directly, $15$ require
the global \texttt{Ec} charge-conjugation convention, and $8$ require
documented row-specific packaging translations. Mutation tests further show
that the comparison detects deliberate changes in signs, coefficients and
index structures (\Cref{sec:validation-sensitivity}).

In this sense, the validation framework is itself part of the result. It gives
a reproducible way to decide whether independently generated tensor-valued
Feynman rules represent the same interaction.

% --------------------------------------------------------------
\section{Lessons from the implementation}
\label{sec:conclusion-hard}
% --------------------------------------------------------------

The main difficulty was not the algebra itself, but the handling of indices.
When a field carries several indices of the same type, an incorrect assignment
can still produce a well-formed expression while contracting the wrong
objects. For this reason, \textsc{FeynPy} avoids guessing ambiguous
assignments and requires them to be resolved explicitly.

Fermionic expressions introduced a second difficulty. Since the underlying
symbolic multiplication is commutative, the information needed to recover
fermion signs must be preserved explicitly throughout the derivation.
Charge-conjugated structures add another layer: two tools can represent the
same physical interaction in different ways even when both are correct. The
validation therefore separates genuine disagreements from differences in
representation rather than forcing every expression into agreement through
ad hoc transformations.

% --------------------------------------------------------------
\section{Limitations and future work}
\label{sec:conclusion-limitations}
\label{sec:conclusion-future}
% --------------------------------------------------------------

The present implementation is limited to tree-level Feynman rules. It does
not include loop calculations, counterterms or renormalisation.

It also produces symbolic vertices rather than external model files. A UFO
interface is therefore the most direct practical extension. Much of the
necessary information is already available in structured form; the remaining
step is to translate particles, parameters and tensor structures into the UFO
representation. This would allow \textsc{FeynPy} models to be used directly
with tools such as \textsc{MadGraph5\_aMC@NLO}~\cite{Alwall:2014hca}.

A second important extension concerns effective field theory. The current
integration-by-parts implementation is restricted to unlabelled scalar
species. Extending it to general fields and adding equation-of-motion
relations would make it possible to reduce redundant operator sets, for
example from a Green basis to the Warsaw basis.

Other limitations are more practical: restriction files are not yet supported,
general open multi-fermion tensor structures are less complete than closed
Dirac bilinears, and the validation compares interaction vertices rather than
parameter-card semantics or numerical metadata.

Gauge fixing could also be extended to the full broken $R_\xi$ gauge, where
gauge bosons and Goldstone fields mix in the gauge-fixing functions.

Finally, the validation can be expanded beyond the three models considered
here. Models with extended gauge sectors, supersymmetry or higher-dimensional
SMEFT operators would provide useful tests of different parts of the
implementation.

% --------------------------------------------------------------
\section{Closing remark}
\label{sec:conclusion-closing}
% --------------------------------------------------------------

Deriving Feynman rules is, in principle, a mechanical procedure. For a theory
as large as SMEFT, however, performing and checking that procedure by hand is
no longer realistic. The goal of this work was therefore not to invent a new
derivation algorithm, but to build the structure around the standard one
carefully enough that it can be applied to realistic theories and its output
can be trusted.

\textsc{FeynPy} combines a compact Python model description, generic gauge
and tensor manipulation, and a canonical validation procedure in a single
framework. Its agreement with the \textsc{FeynRules} reference vertices
across three complete models provides evidence that this approach works. The
result is not yet a replacement for the full \textsc{FeynRules} ecosystem,
but it shows that reliable Feynman-rule derivation and validation can be
carried out natively within Python.
